\documentclass[10pt, handout]{beamer}

%\usepackage[backend=bibtex,firstinits=true,style=verbose-inote,citestyle=authortitle]{biblatex}
\usepackage{bm}
\usepackage{graphicx}
\usepackage{subcaption}
\usepackage{amsmath}
\usepackage{amsfonts}
\usepackage{makecell}
\usepackage{filecontents}
\usepackage{biblatex}
% \newcommand{\expect}[2][]{
\ifthenelse{\equal{#1}{}}{
\mathbb{E}\left[#2\right]
}{
\underset{#1}{\mathbb{E}}\left[#2\right]
}}

\newcommand{\cov}[2][]{
\ifthenelse{\equal{#1}{}}{
\text{Cov}\left[#2\right]
}{
\underset{#1}{\text{Cov}}\left[#2\right]
}}


\newcommand{\var}[2][]{
\ifthenelse{\equal{#1}{}}{
\text{Var}[#2]
}{
\underset{#1}{\text{Var}}[#2]
}}

\newcommand{\loss}[2][]{
\ifthenelse{\equal{#1}{}}{
\mathcal{L}(#2)
}{
\mathcal{L}_{#1}(#2)
}}

\newcommand{\kl}[2]{
\text{D}_\text{KL}[#1 \parallel #2]
}

\newcommand{\R}{\mathbb{R}}
%\newcommand{\Prob}{\mathbb{P}}

\newcommand{\1}[1]{\mathds{1}\{#1\}}


%\usecolortheme{dolphin}
\setbeamertemplate{navigation symbols}{}
\setbeamertemplate{section in toc}{\inserttocsectionnumber.~\inserttocsection}


\addbibresource{references.bib}


\title{Model Compression}
%\subtitle{}
\author{Ivan Skorokhodov}
%\date{}
%\logo{\includegraphics[height=1cm]{images/ipavlov-logo.png}}

\newcommand{\citepaper}[1]{\citetitle{#1} by \citeauthor{#1}}

%\graphicspath{{./images}}

%\usetheme{lucid}
\begin{document}

\begin{frame}
    \titlepage
\end{frame}

\begin{frame}{Overview}
    \tableofcontents
\end{frame}

\section{Motivation}
\begin{frame}
    \begin{itemize}
        \item\pause Accelerating inference or training
        \item\pause Reducing memory footprint
        \item\pause Theoretical curiosity
    \end{itemize}
\end{frame}


\section{Pruning}
\begin{frame}{Pruning}
    \begin{itemize}
        \item\pause \textit{Pruning} is removing weights/neurons in a model while preserving the accuracy
        \item\pause It can be done at different stages:
        \begin{itemize}
            \item\pause before training
            \item\pause during training
            \item\pause after training
            \item\pause iteratively train/prune several times
        \end{itemize}
        \item\pause Pruning neurons speeds up a model, but:
        \begin{itemize}
            \item\pause \cite{Rethinking_Pruning} argues that training the pruned model from scratch would give the same performance
            \item\pause So the main value is in optimizing the architecture
        \end{itemize}
        \item\pause Pruning weights (theoretically) reduces the number of FLOPs, but:
        \begin{itemize}
            \item\pause Resulted sparse matrices are not ``sparse enough'' to provide practical benefits (sparse matrix-vector multiplications are usually based on non-parallel computations)
            \item\pause \cite{State_of_sparsity, Rethinking_Pruning} claim that modern SotA weight-pruning algorithms do  not generalize on large datasets
        \end{itemize}
    \end{itemize}
\end{frame}


\subsection{Pruning weights}
\begin{frame}{Pruning weights}
    \pause Simple strategies:
    \begin{itemize}
        \item Apply $L_1$-regularization during training
        \item Iterative Magnitude Pruning (IMP): prune weights
    \end{itemize}
\end{frame}


\begin{frame}{Lottery Ticket Hypothesis (LTH)}
    \begin{itemize}
        \item 
    \end{itemize}
\end{frame}


\begin{frame}{Synaptic Flow \cite{SynFlow} (part 1/k)}
    \begin{itemize}
        \item\pause Pruning algorithms remove weights based on score values associated with each weight
        \item\pause These scores can usually be represented as
    \begin{equation}\label{eq:syn-saliency}
        \mathcal{S}(\theta)=\frac{\partial \mathcal{R}}{\partial \theta} \odot \theta
    \end{equation}
    for some function $R(\theta)$
        \item\pause Authors call this function $S(\theta)$ a \textit{synaptic saliency} (importance) function
    \end{itemize}
\end{frame}


\begin{frame}
\pause Many pruning algorithms have a form similar to \eqref{eq:syn-saliency}:
    \begin{itemize}
        \item\pause Magnitude Pruning:
        \begin{equation}
            \mathcal{R}(\theta) = \frac{1}{2}\| \theta \|_2^2 \Longrightarrow S(\theta_i) = \theta_i^2
        \end{equation}
        \item\pause Skeletonization \cite{Skeletonization}:
        \begin{equation}
            
        \end{equation*}
    \end{itemize}
\end{frame}


\subsection{Pruning neurons}
\begin{frame}{Pruning neurons (\textit{structured pruning})}
    \begin{itemize}
        \item 
        \item 
    \end{itemize}
\end{frame}



\section{Hashing}
\begin{frame}
    pass
\end{frame}

\subsection{Simple hashing}
\begin{frame}
    pass
\end{frame}

\subsection{Multi-hashing}
\begin{frame}
    pass
\end{frame}

\section{Quantization}
\begin{frame}
    pass
\end{frame}

\section{Low-rank decomposition}
\begin{frame}
    pass
\end{frame}

\begin{frame}{}
    
\end{frame}

\section{Other techniques}
\begin{frame}{Knowledge distillation}
    pass
\end{frame}

\begin{frame}{Conditional computation}
    pass
\end{frame}

\begin{frame}{Architectural tricks}
    pass
\end{frame}

\section{Conclusion}
\begin{frame}
    pass
\end{frame}


\end{document}
